Four consecutive days on the surface of Mars will be shown in the
planetarium. There is a Martian base visible on the horizon. During the
Martian daytime the sky will be slightly brightened. The moons of Mars and
the other planets will not be displayed. The local meridian will be
continuously visible on the sky. You are given a map of the sky for epoch
J2000.0, using a polar projection with a linear scale in declination,
covering stars down to about 5th magnitude.

Note: Azimuth is counted from $0^\circ$ to $360^\circ$ starting at S through
W, N, E.

\begin{description}
\item[(H)] Give the areographic (Martian) latitude of the observer
  $\varphi$.
\item[(I)] Give the altitudes of upper ($h_u$) and lower ($h_l$) culmination
  of:
  Pollux ($\beta$ Gem) and Deneb ($\alpha$ Cyg).
\item[(J)] Give the areocentric (Martian) declination of:
  Regulus ($\alpha$ Leo) and Toliman ($\alpha$ Cen).
\item[(K)] Sketch diagrams to illustrate your working in questions (I) and
  (J) above.
  [Two blank diagrams are provided, each showing the zenith, nadir, and the
  S--N horizon line through the observer.]
\item[(L)] On the map of the sky, mark (with a cross) and label (``M'') the
  Martian celestial North Pole.
\item[(M)] Give the azimuth $A$ of the observer as seen from the Martian
  base.
\item[(N)] Estimate the location of the base on Mars, and circle the
  appropriate description:
  \begin{description}
  \item[(a)] near the Equator
  \item[(b)] near the northern Tropic circle
  \item[(c)] near the northern Arctic circle
  \item[(d)] near the North Pole
  \end{description}
\item[(O)] A time axis is given showing the Martian year and the seasons in
  the northern hemisphere (spring, summer, autumn, winter). Mark the date
  represented by the planetarium display on the axis.
\end{description}

% FIGURE NEEDED: polar-projection sky map (same map as Question 1) for the
% marking task L, and the blank season time axis for task O
